\documentclass{nbeamer}
\usepackage{nbeamer-code}
\nbsetup{grid=true, boxes=true}
\title{nbeamer 设计说明}
\subtitle{网格布局}
\author{XIAN}
\institute{纳米树研究室 Namishu Research Lab}
\date{\today}

\begin{document}
\begin{frame}
  \titlepage
\end{frame}

\begin{frame}
  \title{目录}
  \usebeamertemplate{contents page}
\end{frame}

\section{版面}
\begin{frame}[fragile]{网格}
  整张画布横向 32 格，纵向 18 格。每格都是正方形，边长等于页宽的 $1/32$。在文稿中，用 \incode{\nbgrid} 表示这个长度。

  {\centering
    \begin{tikzpicture}[x=.5cm,y=.5cm]
      \draw[step=.5cm,gray] (0,0) grid (8,2);
      \draw[thick] (3,0) rectangle (4,1);
      \draw[<->] (3,-.45) -- node[below]{一格} (4,-.45);
    \end{tikzpicture}\par}

  数网格可以读出边距和区域宽度。文字另有基线，同一行的字符沿基线排列。
\end{frame}

\begin{frame}{页边距}
  从页面左边缘数到辅助框，有 2 格留白；右边也留出 2 格。页面横向共 32 格，中间可用的宽度是 28 格。

  标题区和内容区共用左右边界。调整任意一侧的页边距，两个区域在该侧的边线都会随之移动。
  \vfill
  内容区底边 \hfill 距页底 2 格
\end{frame}

\section{页面}
\begin{frame}{封面}
  \begin{columns}
    \begin{column}{13\nbgrid}
      \setlength{\lineskip}{0pt}% Align the tall picture with the column top.
      \begin{figure}
        \begin{tikzpicture}[
          x={\dimexpr(\linewidth-\pgflinewidth)/32\relax},
          y={\dimexpr(\pgflinewidth-\linewidth)/32\relax},
          every node/.style={inner sep=1pt,font=\tiny}]
          \draw (0,0) rectangle (32,18);
          \draw[gray] (3,3) rectangle (29,7) (3,7) rectangle (29,9)
            (3,11) rectangle (29,13) (3,13) rectangle (29,14)
            (3,15) rectangle (29,16);
          \node[anchor=west] at (3,5) {nbeamer 设计说明};
          \node[anchor=west] at (3,8) {网格布局};
          \node[anchor=west] at (3,12) {作者};
          \node[anchor=west] at (3,13.5) {机构};
          \node[anchor=east] at (29,15.5) {日期};
        \end{tikzpicture}
      \end{figure}
    \end{column}
    \hfill
    \begin{column}{13\nbgrid}
      \justifying
      封面的左右边界各比普通页向内收 1 格，默认宽 26 格。

      标题和副标题共用左边线。署名放在下方，作者和机构各占一行；日期沿容器右边线对齐。封面标题和正文页标题均限一行，超宽时截断并在编译日志中给出警告。
    \end{column}
  \end{columns}
\end{frame}

\begin{frame}{导航页}
  \begin{columns}
    \begin{column}{13\nbgrid}
      \justifying
      \textbf{目录}

      左侧放页名，右侧从上到下排列章节。每项占 2 格高，最多五项。本稿有三个分类，占前三个位置。
    \end{column}
    \hfill
    \begin{column}{13\nbgrid}
      \justifying
      \textbf{章节页}

      需要单独引出一章时，可在同样宽的容器中只放章节名。它与目录页名共用左边线和标题行的位置。
    \end{column}
  \end{columns}
\end{frame}

\begin{frame}{正文页}
  \begin{columns}
    \begin{column}{13\nbgrid}
      \setlength{\lineskip}{0pt}
      \begin{figure}
        \begin{tikzpicture}[
          x={\dimexpr(\linewidth-\pgflinewidth)/32\relax},
          y={\dimexpr(\pgflinewidth-\linewidth)/32\relax},
          every node/.style={font=\scriptsize}]
          \draw (0,0) rectangle (32,18);
          \draw (2,1) rectangle (30,3);
          \node[font=\tiny] at (16,2) {标题区};
          \draw[fill=black!6] (2,4) rectangle (30,16);
          \node at (16,10) {内容区};
        \end{tikzpicture}
      \end{figure}
      \justifying
      \textbf{有标题页。}标题区位于上方，内容区从距页顶 4 格处开始，距页底留 2 格。
    \end{column}
    \hfill
    \begin{column}{13\nbgrid}
      \setlength{\lineskip}{0pt}
      \begin{figure}
        \begin{tikzpicture}[
          x={\dimexpr(\linewidth-\pgflinewidth)/32\relax},
          y={\dimexpr(\pgflinewidth-\linewidth)/32\relax},
          every node/.style={font=\scriptsize}]
          \draw (0,0) rectangle (32,18);
          \draw[fill=black!6] (2,2) rectangle (30,16);
          \node at (16,10) {内容区};
        \end{tikzpicture}
      \end{figure}
      \justifying
      \textbf{无标题页。}内容区从距页顶 2 格处开始，比左图多出 2 格高度。左右边界和底边与左图相同。
    \end{column}
  \end{columns}
\end{frame}

\section{内容}
\begin{frame}[fragile]{分栏}
  \begin{columns}
    \begin{column}{13\nbgrid}
      \justifying
      左栏从普通文字开始。首行基线位于栏顶下方 $\tfrac34$ 格，汉字和字母沿这条线排列。

      本页两栏各宽 \incode{13\nbgrid}，中间留 2 格。嵌套分栏时，每格仍是页宽的 $1/32$。
    \end{column}
    \hfill
    \begin{column}{13\nbgrid}
      \begin{table}
        \begin{tabular*}{\linewidth}{l@{\extracolsep{\fill}}l}
          \toprule 栏首内容 & 对齐参照 \\
          \midrule 正文 & 首行基线 \\
          图片 & 图像上边缘 \\
          表格 & 表体上边缘 \\ \bottomrule
        \end{tabular*}
        \caption{表格撑满右栏，顶线贴住栏顶。}
      \end{table}
    \end{column}
  \end{columns}
\end{frame}

\begin{frame}{段落}
  第一行和下一行属于同一段。\\
  从上一条基线量到这一条，距离是 \textbf{1 格}。

  这里开始第二段。上一段末行到这一段首行，基线相距 \textbf{2 格}。多出的一格标明段落边界。

  普通字号的正文沿用这个节奏。\emph{强调}和\textbf{粗体}让读者在同一行里辨认重点。
\end{frame}

\begin{frame}{环境}
  \begin{columns}
    \begin{column}{14\nbgrid}
      \begin{block}{内容块}
        标题和首句同行。

        块内另起一段，两条普通文字基线相距 2 格。
      \end{block}
      这里接回正文。列表、块或证明结束后，到下一段普通正文也留出同样的段距。
    \end{column}
    \hfill
    \begin{column}{12\nbgrid}
      \begin{figure}
        \includegraphics[width=\linewidth]{../nbeamer/figures/mushroom-shelter.png}
        \caption{题注随栏宽换行。}
      \end{figure}
    \end{column}
  \end{columns}
  图表和高公式按实际高度展开。后续内容可以从两条网格线之间开始。
\end{frame}
\begin{frame}{图片网格}
  图片等宽排列。相邻两行之间留 1 格，四列布局的列间距为 1 格，其他多列布局的列间距为 2 格。扣除间距后，剩余宽度由各列均分，图片宽度可以是小数格。

  图片网格题注沿排字方向只占一行。超出单元格可用长度时，保留开头并截断末尾，编译继续，日志提示对应题注。普通图表的题注随栏宽换行。
\end{frame}
\end{document}
